Veľké jazykové modely (LLM) dnes vedia plynule odpovedať takmer na čokoľvek. Problém je, že odpovede znejú vierohodne aj vtedy, keď model za sebou nemá žiadny konkrétny zdroj. V právnej doméne, kde každé tvrdenie musí byť oprené o konkrétny paragraf, je takéto správanie diskvalifikujúce.

Bežná odpoveď na halucinácie je RAG (Retrieval Augmented Generation): pred vygenerovaním odpovede sa do promptu vloží relevantný kúsok textu nájdený vo vektorovej databáze. Pre krátke faktické otázky tento prístup stačí. Pre slovenský zákon nie. Zákon nie je súvislý dokument, kde je odpoveď vždy na jednom mieste. Definícia základu dane môže byť v jednom paragrafe, výnimky o desiatky strán ďalej, sadzba v prílohe a zmysel z toho vznikne až po prejdení reťaze krížových odkazov. Vektorové vyhľadávanie tieto skoky nezachytí, lebo sa pozerá iba na sémantickú podobnosť textu, nie na to, čo na koho odkazuje.

Znalostný graf takéto väzby modeluje priamo. Entity sú uzly, odkazy medzi paragrafmi sú hrany, a jazykový model sa namiesto hľadania podobných viet pohybuje po vopred extrahovanej štruktúre dokumentu. Tabuľky a vzorce, ktoré v plochom textovom RAG-u úplne padajú (tabuľky sa rozsypú na nesúrodé riadky, vzorce zostanú v PDF ako prázdne obrázky), sa do grafu dajú vložiť ako samostatné typy uzlov s vlastnými vlastnosťami a zachovať tak ich pôvodný význam.

Práca je členená do štyroch hlavných kapitol. Kapitola \ref{sec:sucasny-stav} \emph{Súčasný stav} zhŕňa, čo sú znalostné grafy, prečo sa hodia ako uzemnenie pre LLM, ako vyzerali klasické metódy extrakcie entít a vzťahov pred príchodom jazykových modelov, a aké existujúce prístupy už na automatickú tvorbu KG z textu existujú. Záver kapitoly otvára medze tých istých prístupov v slovenskom právnom kontexte. Kapitola \ref{sec:navrh} \emph{Návrh riešenia} popisuje samotnú architektúru pipeline od predspracovania PDF cez trojkrokovú extrakciu (objav typov, zjednotenie schémy, schémou riadená extrakcia) až po vkladanie uzlov a hrán do Neo4j s vektorovými indexmi. Kapitola \ref{sec:implementacia} \emph{Implementácia} sa venuje technologickému stacku, riešeniu konkrétnych problémov ako rate-limity API a štruktúrovaný výstup z modelu, chunkovacej stratégii podľa štruktúry zákona a prompt engineeringu pre slovenčinu. Posledná kapitola \ref{sec:experimenty} \emph{Experimenty a vyhodnotenie} prináša kvantitatívne porovnanie modelov a režimov promptovania voči ručne anotovanej referencii, kvalitatívnu analýzu typických chýb a porovnanie odpovedí GraphRAG-u s klasickým vektorovým RAG-om nad rovnakým testovacím zákonom o DPH.
